% chapters/06_fazit_und_ausblick.tex
%
% Erstentwurf vom 27.10.2026 (KI-gestützt erstellt, siehe Erklärung).
% 6.1 folgt den Forschungsfragen FF1 bis FF3 aus Abschnitt 1.2; die
% Antworten enthalten \PH{...}-Platzhalter, die nach Kapitel 5 zu füllen
% sind. 6.2 und 6.3 sind aus Kapitel 4 und 5 ableitbar und ausformuliert.
\providecommand{\PH}[1]{\textcolor{red}{\textbf{[#1]}}}

\chapter{Fazit und Ausblick}
\label{chap:fazit}

Diese Arbeit hat zwei Klassen vortrainierter Sprachmodelle für
feingranulare Named Entity Recognition in einer gemeinsamen Pipeline
verglichen: DeBERTa-v3-base und -large mit vollständigem Fine-Tuning
sowie Qwen3.5-0.8B, -4B und -27B im Zero-Shot-Regime und nach einer
Anpassung mit QLoRA.
Datenbasis waren die offiziellen englischen Splits von MultiNERD mit
15 Entitätstypen; bewertet wurde mit strikten Entity-Level-Metriken,
ergänzt um Trainingsdauer, Inferenzlatenz, GPU-Spitzenspeicher und die
Formtreue der generierten Ausgaben.
Abschnitt~\ref{sec:fazit-zusammenfassung} beantwortet die drei
Forschungsfragen aus Abschnitt~\ref{sec:problem_ziel},
Abschnitt~\ref{sec:limitationen} benennt die Grenzen der Untersuchung,
Abschnitt~\ref{sec:ausblick} skizziert weiterführende Arbeiten.

\section{Zusammenfassung der zentralen Ergebnisse}
\label{sec:fazit-zusammenfassung}

% TODO: Zahlen aus tab:haupt, tab:effizienz und tab:formtreue-erg
% übernehmen; die Richtung der Aussagen an Abschnitt 5.7 angleichen.
% Keine neuen Ergebnisse einführen, die nicht in Kapitel 5 stehen.
% Nach dem Füllen jeden FF-Absatz auf sechs bis acht Sätze kürzen; Typ-
% und Formtreue-Details bleiben in Kapitel 5.

\paragraph{FF1: Erkennungsleistung und Formtreue.}
Auf dem englischen MultiNERD-Testsplit erreicht das beste LLM,
Qwen3.5-27B im QLoRA-Regime, einen Micro-F$_1$ von \PH{xx,x} Punkten,
der beste Encoder, DeBERTa-v3-large, einen von \PH{xx,x}.
Der Abstand von \PH{x,x} Punkten ist \PH{größer / kleiner} als die
Streuung über die Seeds; \PH{das LLM / der Encoder / keine der beiden
Klassen} ist damit unterscheidbar besser.
Die kleineren LLMs erreichen nach der Anpassung \PH{xx,x} (4B) und
\PH{xx,x} Punkte (0.8B) und liegen damit \PH{über / unter / auf dem
Niveau von} DeBERTa-v3-base mit \PH{xx,x}.
Ohne Training fallen die LLMs auf \PH{xx} bis \PH{xx} Punkte zurück;
\PH{die Anpassung mit QLoRA ist damit der mit Abstand wirksamste Faktor
im Vergleich, nicht die Modellklasse und nicht die Modellgröße}.
Bei den drei häufigen Typen PER, ORG und LOC unterscheiden sich die
trainierten Konfigurationen um \PH{höchstens x,x} Punkte; die
Unterschiede konzentrieren sich auf \PH{die verrauschten Typen ANIM, FOOD
und PLANT sowie auf DIS, EVE und MEDIA}.
Die Formtreue ist nach dem Fine-Tuning \PH{hoch}: Weniger als
\PH{x}\,\% der Antworten sind nicht parsebar, unbekannte Typen und
Text-Abweichungen treten \PH{praktisch nicht mehr} auf, und die
verbleibenden Parse-Fehler gehen \PH{nahezu vollständig} auf das Limit
von 256 Ausgabetokens zurück.
Im Zero-Shot-Regime ist die Formtreue \PH{deutlich schwächer} und
\PH{beim 0.8B-Modell der begrenzende Faktor}.
Das JSON-Format mit expliziten Wortindizes hat sich als \PH{tragfähig}
erwiesen.

\paragraph{FF2: Modellgröße und Ressourcen.}
Innerhalb der Encoder-Familie bringt der Schritt von 184 auf 434~Mio.
Parameter \PH{x,x} Punkte, innerhalb der LLM-Familie im QLoRA-Regime der
Schritt von 0{,}75 auf 4{,}21~Mrd. \PH{x,x} Punkte und der von 4{,}21 auf
26{,}90~Mrd. weitere \PH{x,x} Punkte.
\PH{Die Zugewinne nehmen mit der Größe ab / Der Schritt auf 27B ist nicht
unterscheidbar.}
Dem stehen Ressourcen gegenüber, die mit der Größe \PH{nahezu linear}
wachsen: Die Inferenzlatenz je Satz steigt von \PH{8,7}~ms
(DeBERTa-v3-base) über \PH{xx,x}~ms (DeBERTa-v3-large) und \PH{xxx}~ms
(Qwen3.5-0.8B) auf \PH{xxxx}~ms (Qwen3.5-27B), der GPU-Spitzenspeicher
der Inferenz von \PH{x,x} auf \PH{xx,x}~GB, die Trainingsdauer von
\PH{x,x} auf mindestens \PH{xx} Stunden.
Im Zero-Shot-Regime ist der Größeneffekt \PH{größer}: \PH{Erst ab 4B /
27B} erzeugen die Modelle in nennenswertem Umfang gültige Ausgaben.
Über die Familien hinweg erklärt die Parameterzahl die Leistung nicht:
DeBERTa-v3-base erreicht mit 184~Mio. Parametern \PH{mehr / ebenso viel}
wie Qwen3.5-0.8B mit dem Vierfachen.

\paragraph{FF3: Empfehlungen.}
\PH{Unter Latenz- oder Speicherbeschränkungen ist DeBERTa-v3 die
effizientere Wahl}: \PH{DeBERTa-v3-large} erreicht \PH{xx,x} Punkte bei \PH{xx}~ms je
Satz und \PH{x,x}~GB GPU-Speicher und verliert gegenüber dem besten LLM
\PH{höchstens x,x} Punkte, bei einem \PH{xx}-fach geringeren
Latenz- und einem \PH{x}-fach geringeren Speicherbedarf.
Ein LLM rechtfertigt den Mehraufwand nur, wenn \PH{der Leistungsvorsprung
von x,x Punkten für die Anwendung zählt und} Latenzen von \PH{rund x}
Sekunden je Satz sowie \PH{xx}~GB GPU-Speicher akzeptabel sind; dann ist
Qwen3.5-27B mit QLoRA die beste Konfiguration, \PH{während das 4B-Modell
einen Mittelweg bietet / während das 4B-Modell von DeBERTa-v3-large
dominiert wird}.
Das Zero-Shot-Regime ist bei 15 Typen \PH{nicht konkurrenzfähig} und nur
dann sinnvoll, wenn keine annotierten Daten vorliegen.
Diese Empfehlungen gelten für die Verarbeitung einzelner Sätze auf den
verwendeten GPU-Typen; für Offline-Verarbeitung mit Batching verschieben
sich die Faktoren zugunsten der LLMs, ohne dass dies hier gemessen wurde.

\section{Limitationen der Arbeit}
\label{sec:limitationen}

Die Aussagekraft der Ergebnisse ist durch das Versuchsdesign und die
Datenbasis in mehreren Punkten begrenzt; Tabelle~\ref{tab:abweichungen}
fasst die bekannten Abweichungen vom idealisierten Versuchsdesign
zusammen.

\paragraph{Datenbasis.}
Alle Werte messen die Übereinstimmung mit einer automatisch erzeugten
Referenz.
Die MultiNERD-Splits sind per Distant Supervision annotiert, und die
Originalarbeit beziffert die Qualität der Annotation für ANIM, BIO, FOOD
und PLANT mit F$_1$-Werten zwischen 0{,}52 und 0{,}71 gegenüber einer
manuellen Prüfung \cite{tedeschi2022multinerd}.
Ein Teil der gemessenen Fehler bei diesen Typen kann daher auf die
Referenz zurückgehen, und ein Modell, das die Silberannotation
\enquote{zu gut} reproduziert, lernt auch deren Fehler.
Die Untersuchung beschränkt sich zudem auf einen Datensatz, eine Sprache
und ein Genre; Wikipedia- und Wikinews-Texte sind mit hoher
Wahrscheinlichkeit Teil der Vortrainingsdaten der LLMs, und eine
Kontamination wurde nicht empirisch geprüft
(Abschnitt~\ref{sec:research_gap}).
Der CrossNER-Transfertest mildert die Beschränkung auf ein Genre nur
teilweise, da er auf drei Grobtypen reduziert ist und 21 seiner Sätze
mit MultiNERD-Trainingssätzen übereinstimmen.

\paragraph{Vergleichsbedingungen.}
Die Arbeit vergleicht Konfigurationen, nicht isolierte Faktoren.
Die Encoder wurden vollständig feinjustiert, die LLMs nur über Adapter
mit 0{,}5 bis 0{,}8\,\% trainierbarer Parameter; das entspricht der
praxisüblichen Anpassung je Modellklasse, bedeutet aber, dass ein
vollständiges Fine-Tuning der LLMs oder höhere LoRA-Ränge zu anderen
Ergebnissen führen könnten.
Die Adapter decken wegen der hybriden Architektur nur die
Feed-Forward-Netze aller Schichten und die Attention jeder vierten
Schicht ab, und die QLoRA-Umsetzung weicht in zwei Punkten vom
Standardrezept ab (Abschnitt~\ref{sec:lora-konfiguration}).
Hyperparameter wurden je Modellgröße nach üblichen Werten gesetzt und
nicht gesucht; für keine Konfiguration liegt daher ein Nachweis vor,
dass sie ihr Optimum erreicht hat.
Die Checkpoint-Auswahl ist asymmetrisch: Die Encoder wurden auf dem
vollständigen Validierungssplit ausgewählt, die LLM-Adapter anhand von
200 Validierungssätzen.
Der Trainingsverlust der LLMs umfasst den konstanten Prompt, sodass ein
Teil des Lernsignals auf dessen Reproduktion entfällt.

\paragraph{Zero-Shot-Regime.}
Die Zero-Shot-Ergebnisse beruhen auf je einem Lauf mit Sampling bei
Temperatur 1{,}0 und sind damit nicht deterministisch; eine Streuung
über Wiederholungen wurde nicht gemessen.
Prompt und Decoding-Einstellungen wurden nach Beobachtungen auf dem
Testsplit festgelegt, und die Zero-Shot-Läufe verwenden einen anderen
Prompt, eine andere Präzision und ein anderes Decoding als die
QLoRA-Läufe.
Die Zero-Shot-Werte sind deshalb als optimistische Einzelwerte für genau
dieses Format zu lesen; mit Demonstrationsbeispielen oder Constrained
Decoding wären höhere Werte zu erwarten.

\paragraph{Effizienzmessung.}
Die Konfigurationen liefen auf drei GPU-Typen; Latenz- und
Speicherwerte verschiedener Konfigurationen enthalten daher
Hardwareeffekte, die den gemessenen Unterschied zwischen A100 und H100
beim Encoder (17{,}6 gegenüber 8{,}7~ms) erreichen können.
Gemessen wurde die Latenz bei Batchgröße~1 und einem Satz je Anfrage;
der mit Batching erreichbare Durchsatz, der Speicherbedarf während des
Trainings sowie Energieverbrauch und Kosten wurden nicht erfasst.
Die Trainingsdauer des 27B-Modells ist nur als Untergrenze bekannt.
Das Limit von 256 Ausgabetokens führt bei langen Sätzen zu
abgeschnittenen, nicht parsebaren Antworten und unterschätzt Recall und
F$_1$ der LLMs\PH{, auf CrossNER stärker als auf MultiNERD}.

\paragraph{Statistische Absicherung.}
Jede trainierte Konfiguration wurde mit drei Seeds trainiert; auf
inferenzstatistische Tests wurde verzichtet.
Unterschiede in der Größenordnung der Standardabweichungen, darunter
\PH{der Vergleich der beiden Encoder und der Schritt von 4B auf 27B},
lassen sich mit diesem Umfang nicht belegen.
Der Testsplit wurde außerdem bereits in den explorativen Vorläufen
ausgewertet; für die trainierten Konfigurationen blieben die
Hyperparameter danach bis auf die Epochenzahl des 27B-Modells
unverändert, eine strenge Trennung von Entwicklung und Test ist aber
nicht gegeben.

\section{Ausblick}
\label{sec:ausblick}

Aus den Ergebnissen und ihren Grenzen ergeben sich mehrere
Anschlussfragen.

\paragraph{Gold-Referenz.}
Die wichtigste Ergänzung wäre eine manuell geprüfte Teststichprobe,
etwa im Umfang der 1\,000 Sätze, die \textcite{tedeschi2022multinerd}
für ihre Qualitätsmessung verwendet haben.
Damit ließe sich trennen, ob die Unterschiede zwischen den Modellklassen
bei ANIM, FOOD und PLANT Modellfehler oder Annotationsfehler sind, und
ob ein Modell, das von der Silberreferenz abweicht, in Wahrheit die
bessere Annotation liefert.

\paragraph{Anpassungsverfahren.}
Der Vergleich zwischen vollständigem Fine-Tuning und Adaptern ließe sich
auf der LLM-Seite schließen, indem das 0.8B-Modell vollständig
feinjustiert wird; umgekehrt könnten die Encoder mit LoRA angepasst
werden, wie es \textcite{hu2021lora} für RoBERTa und DeBERTa zeigen.
Auf der Adapterseite bleiben höhere Ränge, Adapter in den
Gated-DeltaNet-Schichten, ein auf die Antwort beschränkter Verlust und
der Verzicht auf Quantisierung offen.
Eine dritte Klasse, Encoder-Generalisten wie GLiNER
\cite{zaratiana2024gliner} oder mit LLM-Annotationen vortrainierte
Encoder wie NuNER \cite{bogdanov2024nuner}, könnte den Vergleich um
Modelle ergänzen, die Vorwissen aus LLMs mit der Effizienz von Encodern
verbinden.

\paragraph{Zero-Shot und Ausgabeformat.}
Für das Zero-Shot-Regime wären mehrere Samples je Satz oder ein
Greedy-Kontrolllauf bei allen Größen nötig, um die Streuung zu
beziffern; Few-Shot-Prompting mit Demonstrationsbeispielen und
Constrained Decoding, etwa das LogitMatch-Verfahren von
\textcite{semin2026spanLabeling}, sind die naheliegenden Erweiterungen.
Für das Ausgabeformat steht ein Vergleich der Wortindex-Variante mit
Inline- und Matching-Formaten unter gleichen Trainingsbedingungen aus,
der zugleich den Einfluss der Ausgabelänge auf die Latenz beziffern
würde \cite{zhan2026generativeNER,lu2024padellm}.
Ein höheres Ausgabelimit würde die Abbrüche bei langen Sätzen
beseitigen.

\paragraph{Effizienz.}
Eine Messung auf einheitlicher Hardware, mit Batching und mit
Energieverbrauch würde die Empfehlungen aus FF3 auf
Offline-Verarbeitung ausweiten.
Für die Inferenz der LLMs bieten sich zudem spezialisierte
Inferenzsysteme und das Zusammenführen der Adapter mit dem Basismodell
an, die hier nicht verwendet wurden.

\paragraph{Übertragbarkeit.}
MultiNERD umfasst zehn Sprachen mit demselben Typinventar.
Die Pipeline dieser Arbeit ließe sich ohne Änderung des Formats auf die
übrigen Sprachen anwenden und würde zeigen, ob der hier für Englisch
gemessene Abstand zwischen Encodern und LLMs in Sprachen mit weniger
Vortrainingsdaten anders ausfällt.
Der CrossNER-Versuch legt nahe, den Domänenwechsel mit dem vollständigen
Typinventar und mit wenigen Zieldomänenbeispielen zu untersuchen, wie es
\textcite{ashok2023promptner} und \textcite{sainz2024gollie} für
LLM-Ansätze tun.
